\subsection{When the joint fibre characters are exactly quadratic}

\begin{statement}{A.80}{Proved}
{A squareclass criterion for the complete common character group}
Retain the finite etale algebra and the actual splitting-field
action of A.57, with $4\mid n$ and $D_1,\ldots,D_k\in\Q^\times$.
Put
\[
 K=\Q(\sqrt{D_1},\ldots,\sqrt{D_k}),\qquad
 V=\langle[D_1],\ldots,[D_k]\rangle
                      \subset\Q^\times/\Q^{\times2}.
 \tag{A.80.1}
\]
In the range $[-1]\notin V$, the largest abelian field contained in
every field factor of
\[
             \Q[Z_1^{\pm1},\ldots,Z_k^{\pm1}]/(Z_i^n-D_i)
\]
is exactly $K$.
Equivalently, all its common finite-order characters are the
quadratic characters in the span of the $D_i$.
Thus an additional joint character requires
$[-1]\in V$.

For an old cover $T-c_i=a_iX_i^n$ and a new rational puncture $c$,
set $D_i=(c-c_i)/a_i$.  In the range $[-1]\notin V$, every
cyclic Brauer class $(T-c,\chi)$ supplied by a common fibre
character is the restriction of an explicit constant class
in $\operatorname{Br}(\Q)[2]$.
In particular this applies whenever all the $D_i$ are positive.
\end{statement}
\begin{proof}
We specify the quotient maps, including the difference between
individual stabilizers and the subgroup they generate.
Use the fields $L/F/\Q$, the group $G=\operatorname{Gal}(L/\Q)$,
and $\Gamma=\operatorname{Gal}(F/\Q)$ of A.57.  Write
$\pi:G\twoheadrightarrow\Gamma$ for restriction and $\Omega$
for the finite set of root tuples of the fibre algebra.
For $\omega=(z_1,\ldots,z_k)\in\Omega$ the actual maps are
\[
 G_\omega\lhook\joinrel\longrightarrow G
       \xrightarrow{\ \pi\ }\Gamma
       \xrightarrow{\ \rho_K\ }\operatorname{Gal}(K/\Q),
 \qquad
 \rho_K(\sigma)(\sqrt{D_i})
             =(-1)^{\eta_i(\sigma)}\sqrt{D_i},
 \tag{A.80.2}
\]
where, in the notation $D_i=(-1)^{\delta_i}\prod_P P^{w_{iP}}$,
\[
 \eta_i(\sigma)=
       \frac{u_\sigma-1}{2}\delta_i
             +\sum_Pw_{iP}\epsilon_P(\sigma)\pmod2.
 \tag{A.80.3}
\]
The root $z_i^{n/2}$ is one of the two square roots of $D_i$.
An element fixing $\omega$ therefore fixes every $\sqrt{D_i}$.
Consequently every composite in (A.80.2) is zero.  In particular,
if $H$ denotes the subgroup generated by the images
$\pi(G_\omega)$, then
\[
 H\subset\Gamma_K:=\ker\rho_K=\operatorname{Gal}(F/K),
 \qquad K\subset F^H=B.
 \tag{A.80.4}
\]
All of these are maps of finite Galois groups; the field inclusion
also follows separately inside each actual root-field factor.

We prove the reverse inclusion using the explicit stabilizer
test, without replacing the union of stabilizers by a subgroup
before taking its generated subgroup.  Suppose
$\sigma\in\Gamma_K$ and $u_\sigma\equiv3\pmod4$.
Write
\[
 d=\gcd(n,u_\sigma-1)=2d_o,
                  \qquad d_o\text{ odd}.
 \tag{A.80.5}
\]
Choose each $k_P\in\Z/n$ with the prescribed parity
$k_P\equiv\epsilon_P(\sigma)\pmod2$ and
$k_P\equiv0\pmod {d_o}$.  The Chinese remainder theorem
gives these choices, since $d_o$ is odd and divides $n$.
For every $i$, the integer
\[
       \frac{u_\sigma-1}{2}\delta_i+\sum_Pw_{iP}k_P
 \tag{A.80.6}
\]
is divisible by $d_o$: here $d_o\mid(u_\sigma-1)/2$.
Its reduction modulo two is $\eta_i(\sigma)=0$, because
$\sigma$ fixes $K$.  It is therefore divisible by $d$.
The exact test (A.57.10) supplies $e_i\in\Z/n$ for which
\[
 (u_\sigma-1)(\delta_i+2e_i)
                       +2\sum_Pw_{iP}k_P=0\pmod {2n}.
 \tag{A.80.7}
\]
The lift $(\sigma,(k_P)_P)\in G$ of (A.57.8) fixes the
root tuple (A.57.9) with these $e_i$.  Thus $\sigma$ belongs
to the image of one actual stabilizer, and hence to $H$.

Quadratic Kummer theory identifies the quadratic subfields of
$K$ with the nonzero elements of $V$.  Therefore
$[-1]\notin V$ is precisely $i\notin K$.
Since $i\in F$, there is an element
$\tau\in\Gamma_K$ acting nontrivially on $i$; this says
$u_\tau\equiv3\pmod4$.
The preceding paragraph proves $\tau\in H$.
For arbitrary $\sigma\in\Gamma_K$, either
$u_\sigma\equiv3\pmod4$, in which case $\sigma\in H$,
or $u_\sigma\equiv1\pmod4$, in which case
$u_{\sigma\tau}\equiv3\pmod4$ and $\sigma\tau\in H$.
Then $\sigma=(\sigma\tau)\tau^{-1}\in H$ as well.
We have proved
\[
                H=\Gamma_K,\qquad B=F^H=K.
 \tag{A.80.8}
\]
Taking the character groups of this exact quotient gives the
claimed complete common character group.

For the Brauer assertion let
$\chi=\sum_i t_i\chi_{D_i}$, $t_i\in\{0,1\}$, be any
such character.  On the old curve put
$z=(T-c)/a_i$; then $z+D_i=X_i^n$.
Because $n$ is even, the mod-two Steinberg relation is
\[
 \{z/X_i^n,D_i/X_i^n\}_2=0,
       \qquad
 \{T-c,D_i\}_2=\{a_i,D_i\}_2.
 \tag{A.80.9}
\]
The first equality has two arguments adding to one; bilinearity
and the fact that $X_i^n$ is a square prove the second equality.
Under the quaternion-symbol map to the Brauer group it gives
the explicit identity
\[
          (T-c,\chi)=\sum_i t_i(a_i,D_i)
                             \quad\text{in }\operatorname{Br}\Q(C).
 \tag{A.80.10}
\]
The right side is one global constant class.  At every ordinary
local point avoiding $c$ this is an equality of its actual local
evaluations, and their global invariant sum is zero by reciprocity.
Finally, if all $D_i>0$, every product representing $V$ is
positive, so $[-1]\notin V$.  This proves the last assertion.
\end{proof}

The certificate \path{code/fibre_squareclass_certificate.py}
enumerates the finite Galois groups, the actual stabilizer images,
and their generated subgroups in representative signed and
dependent-squareclass cases.  It checks the reverse inclusion
(A.80.8) exactly and retains examples from the excluded
$[-1]\in V$ range where extra characters really occur.
The proof above supplies the general assertion for every
$4\mid n$ and every finite tuple of rational radicands.
